\documentclass[12pt]{article}
\usepackage[utf8]{inputenc}
\usepackage[T1]{fontenc}
\usepackage{hyperref}
\title{Kunigami Language Resources for K-12 Curriculum Development}
\author{Compiled by Claude Sonnet 5.5 (claude-sonnet-5-5) for Andr\'as Kornai}
\date{}
\begin{document}
\maketitle

\section{Introduction}
Kunigami (ISO 639-3 xug; Glottocode kuni1268) is an endangered Ryukyuan language of the Japonic family. It is spoken in northern Okinawa Island (the Yanbaru area) and on nearby islands. The Nakijin dialect is often taken as representative. My search found little teaching material. The resources listed are mainly reference works.

\section{Primary Educational Resources}
I found no primers, textbooks or school programmes for Kunigami in the search results.

\section{Linguistic and Reference Materials}
\begin{itemize}
\item Kunigami language (Wikipedia). It gives an overview of the language and its dialects. \url{https://en.wikipedia.org/wiki/Kunigami_language}
\item Oshio, Mutsuko (1999), \emph{Okinawa Iejima h\=ogen jiten} (Dictionary of the Dialect of Ie Island, Okinawa). Glottolog treats Iejima as Kunigami. \url{https://glottolog.org/resource/reference/id/471840}
\item Oshio, Mutsuko (2009), \emph{Shinpan Okinawa Iejima h\=ogen jiten}, the revised edition. \url{https://glottolog.org/resource/reference/id/575735}
\end{itemize}

\section{Religious and Community Resources}
I found nothing specific in the search results.

\section{Digital and Audio Resources}
\begin{itemize}
\item Glosbe Kunigami pages. They appeared to be mainly language-pair pages with little vocabulary, so their usefulness is unconfirmed. \url{https://glosbe.com/bib/xug}
\item Wiktionary has some Kunigami entries. One of them cites the JLect Japonic Languages and Dialects Database. I did not find the JLect site itself. \url{https://en.wiktionary.org/wiki/%E3%81%82%E3%82%93%E3%81%A0%E3%83%BC}
\end{itemize}

\section{Curriculum Development Context}
No school is known to teach Kunigami. Speakers are mostly older adults. A curriculum would have to be a revitalisation effort built on the Ie Island dictionary, other dialect documentation and community recordings. Further searching is needed for Nakijin and Yanbaru materials and for local Okinawan dialect-revival programmes.

\end{document}